\documentclass[leqno]{article}
% Shared mark scheme preamble for resources.danieljsmith.org: the pack
% preamble, unchanged, then the scheme's own layer.
\input{\CurrentFilePath/pack_preamble.tex}

% The page: equal side margins, so the table runs the full width of the
% header rule. The footers' right edge follows the text's.
\geometry{left=0.8in, right=0.8in}
\makeatletter
\renewcommand{\djs@footR}[1]{%
  \put(\dimexpr1in+\hoffset+\oddsidemargin+\textwidth\relax,
    -\dimexpr\paperheight-0.32in\relax){\makebox[0pt][r]{#1}}}
\makeatother

\usepackage[export]{adjustbox}
\usepackage{tabularray}
\definecolor{msGrid}{RGB}{176, 172, 164}
\definecolor{msTint}{RGB}{245, 243, 239}

% The header on every page, the front page included, and the booklet label
% under the contents-page title. \djsQuestionPageStyle stays empty.
\newcommand{\djsSchemeHeader}[1]{%
  \djsPackHeader[markscheme]{\textit{Mark scheme} - #1}%
  \renewcommand{\djsPackTitle}{#1}%
  \renewcommand{\djsBookletLabel}{Mark scheme}%
}

% The front page: the contents heading, spaced as tocloft spaces the packs'
% "Questions", then the contents table. A question's name and page number
% both link to its first page.
\newcommand{\msContentsHeading}{%
  \vskip\cftbeforetoctitleskip
  {\noindent\cfttoctitlefont\contentsname\par}%
  \vskip\cftaftertoctitleskip}
\NewTblrEnviron{mscontents}
\SetTblrInner[mscontents]{
  colspec = {Q[wd=26mm, l] X[l] Q[wd=16mm, c] Q[wd=14mm, r]},
  width = \linewidth,
  rowsep = 4pt, colsep = 6pt,
  hlines = {0.4pt, msGrid}, vlines = {0.4pt, msGrid},
  hline{2} = {2.2pt, scarlet},
  row{1} = {font=\small\bfseries},
  row{Z} = {font=\bfseries},
  column{2} = {font=\small},
}
\newcommand{\msContentsQuestion}[1]{\hyperref[q:#1]{Question #1}}
\newcommand{\msContentsPage}[1]{\hyperref[q:#1]{\pageref*{q:#1}}}
\newcommand{\msSep}{\hspace{0.9em}}

% The conventions pages' layout: headings in the documents' own bold, tables
% ruled and padded as the scheme's, roomy lists.
\newcommand{\msHead}[1]{%
  \par\vspace{18pt}%
  {\noindent\bfseries #1\par}%
  \vspace{6pt}}
\NewTblrEnviron{mstblr}
\SetTblrInner[mstblr]{rowsep = 6pt, colsep = 6pt,
  hlines = {0.4pt, msGrid}, vlines = {0.4pt, msGrid}}
\newenvironment{msList}
  {\begin{itemize}[leftmargin=1.2em, itemsep=5pt, topsep=3pt, parsep=0pt]}
  {\end{itemize}}
% An abbreviation's meaning: its expansion, then the rule it stands for in
% small grey type on the next line.
\newcommand{\msAbbr}[2]{#1\newline{\footnotesize\color{djsMuted}#2}}

% A question's first page: the link target and the PDF bookmark, as
% \questionitem gives the packs.
\newcommand{\msQnum}{}
\newcommand{\msQuestion}[1]{%
  \clearpage
  \renewcommand{\msQnum}{#1}%
  \phantomsection\label{q:#1}%
  \addcontentsline{toc}{subsection}{\protect\djsTocEntry{#1}}}

% The scheme's table: Question 16mm, Scheme the rest, Marks 16mm, Notes
% 44mm; every cell ruled in msGrid; the scarlet rule under the head row. A
% long table has no caption or "continued" foot; a continuation page opens
% with a muted line above the repeated head row. The rules across a part's
% rows, the alternative's tint and the total row are the rendered table's.
\NewTblrTheme{ms}{%
  \DefTblrTemplate{firsthead}{default}{}%
  \DefTblrTemplate{middlehead,lasthead}{default}{%
    {\small\itshape\color{djsMuted}Question \msQnum\ continued\par}%
    \vspace{6pt}}%
  \DefTblrTemplate{firstfoot,middlefoot,lastfoot}{default}{}%
}
\NewTblrEnviron{msscheme}
\SetTblrOuter[msscheme]{long, theme=ms, entry=none, label=none}
\SetTblrInner[msscheme]{
  colspec = {Q[wd=16mm, l] X[l] Q[wd=16mm, c] Q[wd=44mm, l]},
  width = \linewidth,
  rowsep = 6pt, colsep = 6pt,
  rowhead = 1,
  vlines = {0.4pt, msGrid},
  hline{1} = {0.4pt, msGrid}, hline{2} = {2.2pt, scarlet},
  hline{Z} = {0.4pt, msGrid},
  column{1} = {font=\bfseries}, column{3} = {font=\bfseries},
  column{4} = {font=\small},
  row{1} = {rowsep=4pt, font=\small\bfseries},
}

% Two marks on a row: each note starts with its code.
\newcommand{\msNote}[2]{\textbf{#1}\enspace #2}
% Their value, the student's own from earlier working, in quotation marks:
% \(\their{1654} - \their{98}\). A straight " would print as a closing quote.
\newcommand{\their}[1]{\text{``}#1\text{''}}
% A full-row figure: the worked solution's own picture, scaled uniformly and
% only down, to the table less the Question column. The picture is drawn
% with the line width it has in the solutions document, at its list depth
% there (the optional argument, 1 by default: a part's working), so one
% sized by \linewidth keeps its size; tabularray sets the cell's \linewidth
% to the cell's measured width, which would shrink it at every pass. The
% widths are measured once, in a solution box at the packs' text width
% (pack_preamble.tex's margins, 0.8in and 1in).
\newlength{\msSpanW}
\makeatletter
\newcommand{\ms@keepLW}[1]{%
  \expandafter\xdef\csname ms@solnLW@#1\endcsname{\the\linewidth}}
\AtBeginDocument{%
  \setlength{\msSpanW}{\dimexpr\textwidth-16mm-28pt\relax}%
  \setbox\z@\vbox{%
    \begin{minipage}{\dimexpr\paperwidth-1.8in\relax}
      \begin{djsSolution}
        \ms@keepLW{0}%
        \begin{enumerate}\item \ms@keepLW{1}%
        \begin{enumerate}\item \ms@keepLW{2}%
        \begin{enumerate}\item \ms@keepLW{3}%
        \end{enumerate}\end{enumerate}\end{enumerate}
      \end{djsSolution}
    \end{minipage}}}
\newcommand{\msFigure}[2][1]{%
  \adjustbox{max width=\msSpanW, max totalheight=0.6\textheight, valign=t}{%
    \hsize=\csname ms@solnLW@#1\endcsname
    \linewidth=\hsize
    #2}}
\makeatother
% A figure mark's feature, in the Scheme column.
\newcommand{\msFeature}[1]{\textit{#1}}
% An alternative's title row.
\newcommand{\msAlternative}[1]{\textit{Alternative:} #1}
% Drafts only: a queried mark's code, the query's why under its note, and a
% flag in its part's place.
\newcommand{\msQueried}[2]{#1~\(\to\) #2?}
\newcommand{\msQuery}[1]{{\itshape\color{djsMuted}Query: #1}}
\newcommand{\msFlag}[2]{\textbf{#1:}\enspace #2}


\djsSchemeHeader{Simple Harmonic Motion}
\djsSitePath{Further Mechanics}{Circular Motion \& Oscillations}
\djsBoards{\djsBoard{Edexcel FM2}{9FM0/4C}\djsBoard{OCR MEI Mech Major}{Y421}}

\djsLastUpdated{6 November 2026}
\begin{document}
\djsSetQuestionCount{10}
\djsAddMarks{1}{15}
\djsAddMarks{2}{14}
\djsAddMarks{3}{14}
\djsAddMarks{4}{10}
\djsAddMarks{5}{11}
\djsAddMarks{6}{9}
\djsAddMarks{7}{15}
\djsAddMarks{8}{14}
\djsAddMarks{9}{17}
\djsAddMarks{10}{20}
\vspace*{6pt}
\djsContentsTitle
\vspace{30pt}
\msContentsHeading
\begin{mscontents}{}
Question & Marks by part & Marks & Page \\
\msContentsQuestion{1} & (a)\,3\msSep (b)\,2\msSep (c)\,2\msSep (d)\,2\msSep (e)\,6 & 15 & \msContentsPage{1} \\
\msContentsQuestion{2} & (a)\,4\msSep (b)\,6\msSep (c)\,4 & 14 & \msContentsPage{2} \\
\msContentsQuestion{3} & (a)\,5\msSep (b)\,3\msSep (c)\,6 & 14 & \msContentsPage{3} \\
\msContentsQuestion{4} & (a)\,6\msSep (b)\,4 & 10 & \msContentsPage{4} \\
\msContentsQuestion{5} & (a)\,5\msSep (b)\,6 & 11 & \msContentsPage{5} \\
\msContentsQuestion{6} & (a)\,3\msSep (b)\,3\msSep (c)\,3 & 9 & \msContentsPage{6} \\
\msContentsQuestion{7} & (a)\,6\msSep (b)\,3\msSep (c)\,6 & 15 & \msContentsPage{7} \\
\msContentsQuestion{8} & (a)\,6\msSep (b)\,5\msSep (c)\,3 & 14 & \msContentsPage{8} \\
\msContentsQuestion{9} & (a)\,2\msSep (b)\,6\msSep (c)\,3\msSep (d)\,6 & 17 & \msContentsPage{9} \\
\msContentsQuestion{10} & (a)\,5\msSep (b)\,3\msSep (c)\,12 & 20 & \msContentsPage{10} \\
\SetCell[c=2]{r} Total & & 139 & \\
\end{mscontents}
\newpage
% The conventions pages, shared by every pack's mark scheme: one source of
% truth, \input after the front page, with no strand or pack logic. The keys
% on the first side, the rules on the second. The scheme layer defines
% \msHead, \msAbbr, msList and mstblr.
\vspace*{6pt}%
{\noindent\cfttoctitlefont Reading the mark scheme\par}%
\vspace{10pt}%
These rules apply throughout, unless a note says otherwise.

\msHead{Marks}
Every mark is worth one, so \textbf{M1 A1} on a row is two marks. Only marks
coded ft follow through.\par
\vspace{10pt}%
\begin{mstblr}{
  colspec = {Q[wd=15mm, l] X[l]},
  width = \linewidth,
  hline{2} = {2.2pt, scarlet},
  column{1} = {font=\bfseries},
  row{1} = {font=\small\bfseries, rowsep=4pt},
}
Code & Meaning \\
M1 & Method: a correct method, applied to this question. A slip doesn't cost it if the method is clear. A formula must be used, not just quoted. \\
dM1 & Dependent method: as M1, but only if the method mark before it was earned, unless its note names another. \\
A1 & Accuracy: a correct answer or step. It needs the method mark it follows. \\
A1ft & As A1, following through from an earlier answer that the note names. \\
A1~AG & The last mark of a part whose answer is printed in the question: complete working, with no errors. \\
B1 & A correct value, fact or statement. It needs no method mark. \\
B1ft & As B1, following through from an earlier answer that the note names. \\
E1 & A written explanation, interpretation or conclusion. Concluding a proof by induction needs every earlier mark. \\
\end{mstblr}

\msHead{Abbreviations}
ft and AG qualify a code in the Marks column. The others appear in the
notes. A grey line gives the rule where the name alone doesn't.\par
\vspace{10pt}%
\begin{mstblr}{
  colspec = {Q[wd=17mm, l] X[l] Q[wd=15mm, l] X[l]},
  width = \linewidth,
  column{1} = {font=\bfseries}, column{3} = {font=\bfseries},
}
ft & \msAbbr{follow through}{Correct work from the student's own earlier wrong value earns the mark.}
  & AG & \msAbbr{answer given}{The answer is printed in the question, so every step to it must be shown.} \\
cao & \msAbbr{correct answer only}{No follow-through from an earlier error.}
  & cso & \msAbbr{correct solution only}{No error anywhere in the part, not just in its answer.} \\
oe & or equivalent & awrt & answers which round to \\
soi & \msAbbr{seen or implied}{Earned if written, or if later working shows it was used.}
  & isw & \msAbbr{ignore subsequent working}{A correct answer, once reached, isn't spoiled by later working.} \\
SC & \msAbbr{special case}{Credit, as the note says, for work outside the main method.}
  & their & the student's own earlier value \\
condone & a slip that doesn't cost the mark & [\;\,] & working the mark doesn't need \\
s.f. & significant figures & d.p. & decimal places \\
\end{mstblr}

\newpage
\msHead{Reading the scheme}
\begin{msList}
\item Each row is a line of working, with the marks it earns beside it.
  Where a row carries two marks, each note starts with its code.
\item \textbf{[\textit{n}]} ends each part with its total.
\item A shaded \textit{Alternative} is a second method for the same total.
  Mark an answer by the method it follows.
\item Any other valid method earns the equivalent marks, unless the question
  names a method or says \textit{hence}.
\item A sketch earns one B1 for each feature named below its figure.
\end{msList}

\msHead{Accuracy and answers}
\begin{msList}
\item Unless a question says otherwise, give answers to 3 significant
  figures, and angles in degrees to 1 decimal place. Angles in radians take
  3 significant figures.
\item A final answer with more figures is accepted if it rounds to the
  scheme's value; one with fewer is not. An answer spoiled by rounding too
  early loses its accuracy mark.
\item Where a question asks for an exact answer, a decimal loses the final
  accuracy mark. Otherwise either is accepted.
\item Wrong or missing units cost nothing unless the question asks for them.
\item An omitted $\mathrm{d}\theta$, $\mathrm{d}x$ or $\mathrm{d}u$ in an
  integral is not penalised where the variable of integration is clear from
  the integrand and the limits. It is penalised where it is genuinely
  ambiguous: a substitution part-way through, two variables in play, or
  limits in the wrong variable.
\item A correct answer with no working earns full marks, unless the question
  needs the working: show that, prove, hence, a method it names, or an exact
  answer. A wrong answer with no working earns nothing. Calculator answers
  are treated the same way.
\item A correct answer, once reached, isn't spoiled by later working (isw).
\end{msList}

\msHead{Value of $\boldsymbol{g}$}
\begin{msList}
\item Take $g = 9.8\,\mathrm{m\,s^{-2}}$ unless a question states a value,
  which is then the only one accepted. Answers that use $g = 9.8$ may be
  given to 2 or 3 significant figures.
\item Where a question doesn't state $g$, an answer worked with $g = 10$
  throughout earns the same marks, keeping any value printed in the
  question. The worked solutions list the final answers with $g = 10$.
\end{msList}


\msQuestion{1}
\begin{msscheme}{
  hline{3-5,7-8,10-11,13-14,16-21} = {2-Z}{0.4pt, msGrid},
  hline{6,9,12,15,22} = {0.4pt, msGrid},
  row{22} = {font=\bfseries},
}
Question & Scheme & Marks & Notes \\
1(a) & {\(\dot{x} = -\tfrac34\pi\sin\left(\tfrac14\pi t\right) + \tfrac{3\sqrt3}{4}\pi\cos\left(\tfrac14\pi t\right)\)\\[4pt] \(\Rightarrow \ddot{x} = \ldots\)} & M1 & {Differentiates their expression for \(x\) twice, each term changing between sine and cosine. Condone sign and coefficient slips.} \\
& {\(\ddot{x} = -\tfrac{3}{16}\pi^2\cos\left(\tfrac14\pi t\right) - \tfrac{3\sqrt3}{16}\pi^2\sin\left(\tfrac14\pi t\right)\)} & A1 & \\
& {\(\ddot{x} = -\tfrac{\pi^2}{16}\left(3\cos\left(\tfrac14\pi t\right) + 3\sqrt3\sin\left(\tfrac14\pi t\right)\right)\)\\[4pt] \(\phantom{\ddot{x}} = -\tfrac{\pi^2}{16}x\), so simple harmonic motion} & A1~AG & {Correct \(\ddot{x}\) shown equal to \(-\frac{\pi^2}{16}x\) without error, with a conclusion of simple harmonic motion. A0 if the conclusion is missing.} \\
\nopagebreak & & [3] & \\
1(b) & {\(T = \frac{2\pi}{\their{\tfrac{\pi}{4}}} = \ldots\)} & M1 & {Uses \(T = \frac{2\pi}{\omega}\) with their \(\omega\), from \(\ddot{x} = -\omega^2 x\) or the argument \(\frac14\pi t\).} \\
& {\(T = 8\,\mathrm{s}\)} & A1 & \\
\nopagebreak & & [2] & \\
1(c) & {\(a = \sqrt{3^2 + \left(3\sqrt3\right)^2} = \ldots\)} & M1 & {Uses \(\sqrt{p^2+q^2}\) with the coefficients \(3\) and \(3\sqrt3\), or finds the greatest value of \(x\) by another valid method. Condone a squaring slip.} \\
& {\(a = 6\,\mathrm{m}\)} & A1 & \\
\nopagebreak & & [2] & \\
1(d) & {\(v_{\max} = \their{6} \times \their{\tfrac{\pi}{4}} = \ldots\)} & M1 & {Uses \(a\omega\) with their amplitude and \(\omega\), or maximises their \(\dot{x}\) by another valid method.} \\
& {\(v_{\max} = \frac{3\pi}{2}\,\mathrm{m\,s^{-1}}\)} & A1 & {oe in terms of \(\pi\).} \\
\nopagebreak & & [2] & \\
1(e) & {\(\left(\tfrac{\pi}{4}\right)^2\left(\their{6}^2 - x^2\right) > \left(\tfrac{3\sqrt3}{4}\pi\right)^2\)\\[4pt] \(\Rightarrow x^2 < \ldots\)} & M1 & {Uses \(v^2 = \omega^2\left(a^2 - x^2\right)\) with their \(a\) and \(\omega\) and the given speed, as an equation or inequality, to find \(x\), or sets their \(\dot{x}\), or \(a\omega\cos\omega t\), equal to the given speed. Condone slips.} \\
& {\(-3 < x < 3\)} & A1 & {Boundary \(x = \pm 3\), or \(|x| < 3\), soi; or, from the speed, a correct cosine or sine value, such as \(\cos\tfrac14\pi t = \frac{\sqrt3}{2}\), or a correct boundary time.} \\
& {\(6\sin\left(\tfrac14\pi t\right) = \their{3} \Rightarrow t = \ldots\)} & M1 & {Sets a sine or cosine expression for \(x\), the original included, equal to their boundary displacement, or one for the speed equal to the given speed, and solves for a time.} \\
& {\(t = \frac23\)} & A1 & {Time \(\frac23\) from \(O\) to \(x = 3\), or a correct equivalent time from another form of \(x\) or from the speed.} \\
& {\(4 \times \their{\tfrac23} = \ldots\)} & M1 & {Uses symmetry to combine their times into the total per oscillation, as four times the time from \(O\) to \(x = 3\) or an equivalent sum or difference.} \\
& {\(\frac83\,\mathrm{s}\)} & A1 & {\(\frac83\) or awrt 2.67.} \\
\nopagebreak & & [6] & \\
\nopagebreak \SetCell[c=2]{r} Total for Question 1 & & 15 & \\
\end{msscheme}

\msQuestion{2}
\begin{msscheme}{
  hline{3-5,7-11,13-16} = {2-Z}{0.4pt, msGrid},
  hline{6,12,17} = {0.4pt, msGrid},
  row{17} = {font=\bfseries},
}
Question & Scheme & Marks & Notes \\
2(a) & {\(\frac{2\pi}{\omega} = \frac{\pi}{2} \Rightarrow \omega = 4\)} & M1 A1 & {\msNote{M1}{Uses \(T = \frac{2\pi}{\omega}\) with \(T = \frac{\pi}{2}\) to find \(\omega\).}} \\
& {\(1.6^2 = \their{4}^2\left(a^2 - 0.3^2\right)\)} & M1 & {Uses \(v^2 = \omega^2\left(a^2 - x^2\right)\) with \(x = 0.3\), \(v = 1.6\) and their \(\omega\). Condone slips in substitution.} \\
& {\(2.56 = 16\left(a^2 - 0.09\right)\)\\[4pt] \(a^2 = 0.25 \Rightarrow a = 0.5\)} & A1~AG & {Correct equation with \(\omega = 4\) solved without error to \(a = 0.5\); combined steps accepted. Verifying with \(a = 0.5\) scores A0.} \\
\nopagebreak & & [4] & \\
2(b) & {\(1^2 = \their{4}^2\left(0.5^2 - x^2\right) \Rightarrow x = \ldots\)} & M1 & {Uses \(v^2 = \omega^2\left(a^2 - x^2\right)\) with \(v = 1\), \(a = 0.5\) and their \(\omega\) to find \(x\), as an equation or inequality; or sets \(a\omega\cos\omega t\) or \(a\omega\sin\omega t\) equal to 1.} \\
& {\(x = \frac{\sqrt{3}}{4}\)} & A1 & {oe, or awrt 0.433. Sign not needed. Or, from the speed, \(\cos 4t = \frac12\), \(\sin 4t = \frac12\) or a correct boundary time.} \\
& {\(0.5\sin 4t = \their{\tfrac{\sqrt{3}}{4}} \Rightarrow t = \ldots\)} & M1 & {Uses \(x = a\sin\omega t\) or \(x = a\cos\omega t\), with their \(\omega\), equated to their boundary displacement, or its speed equated to 1, and solves for \(t\).} \\
& {\(t = \frac{\pi}{12}\)} & A1 & {oe, including \(t = \frac{\pi}{24}\) from an end point with \(x = a\cos\omega t\).} \\
& {\(4 \times \their{\tfrac{\pi}{12}} = \frac{\pi}{3}\,\mathrm{s}\)} & M1 A1 & {\msNote{M1}{Multiplies their time from \(O\) by 4, or uses an equivalent symmetry argument, such as \(4\left(\frac{\pi}{8} - \their{\frac{\pi}{24}}\right)\) from an end point.}} \\
\nopagebreak & & [6] & \\
2(c) & {\(\tfrac12 \times 0.5 \times v^2 = 0.36 \Rightarrow v^2 = 1.44\)\\[4pt] \(1.44 = \their{4}^2\left(0.25 - x^2\right) \Rightarrow x = \ldots\)} & M1 & {Uses \(\frac12 mv^2 = 0.36\) for \(v^2\), then \(v^2 = \omega^2\left(a^2 - x^2\right)\) with \(a = 0.5\) and their \(\omega\) to find \(x\). Condone slips.} \\
& {\(x = 0.4\)} & A1 & {Accept \(\pm 0.4\).} \\
& {\(F = 0.5 \times \their{4}^2 \times \their{0.4}\)} & M1 & {Uses Newton's second law with acceleration of magnitude \(\omega^2 x\), with their \(\omega\) and displacement. Not for the acceleration alone.} \\
& {\(F = 3.2\,\mathrm{N}\)} & A1 & \\
\nopagebreak & & [4] & \\
\nopagebreak \SetCell[c=2]{r} Total for Question 2 & & 14 & \\
\end{msscheme}

\msQuestion{3}
\begin{msscheme}{
  hline{3-5,7-8,10-12} = {2-Z}{0.4pt, msGrid},
  hline{6,9,13} = {0.4pt, msGrid},
  row{13} = {font=\bfseries},
}
Question & Scheme & Marks & Notes \\
3(a) & {\(3.6^2 = \omega^2\left(a^2 - 0.5^2\right)\)\\[4pt] \(1.5^2 = \omega^2\left(a^2 - 1.2^2\right)\)} & M1 A1 & {\msNote{M1}{Uses \(v^2 = \omega^2\left(a^2 - x^2\right)\) at both given positions. Condone a slip in one equation.}\\ \msNote{A1}{Both equations correct.}} \\
& {\(10.71 = 1.19\,\omega^2 \Rightarrow \omega^2 = 9\)} & dM1 A1 & {\msNote{dM1}{Solves their two equations together, by subtracting or dividing, eliminating \(a\) or \(\omega\).}\\ \msNote{A1}{\(\omega^2 = 9\) or \(\omega = 3\); or \(a^2 = 1.69\) where \(\omega\) is eliminated.}} \\
& {\(12.96 = 9\left(a^2 - 0.25\right) \Rightarrow a^2 = 1.69\)\\[4pt] \(a = 1.3\)} & A1~AG & {Correct equations leading without error to \(a = 1.3\); combined steps accepted. Not for substituting \(a = 1.3\) to verify.} \\
\nopagebreak & & [5] & \\
3(b) & {\(\text{Period} = \frac{2\pi}{3}\,\mathrm{s}\)} & B1 & \\
& {\(\omega^2 a = 9 \times 1.3 = 11.7\,\mathrm{m\,s^{-2}}\)} & M1 A1 & {\msNote{M1}{Uses \(\omega^2 a\) with their \(\omega\) and \(a = 1.3\). Not for \(\omega a\).}\\ \msNote{A1}{11.7 cao.}} \\
\nopagebreak & & [3] & \\
3(c) & {\(9|x| \ge 6.3 \Rightarrow |x| \ge 0.7\)} & M1 A1 & {\msNote{M1}{Sets \(\omega^2 x\), with their \(\omega\), equal to \(6.3\), as an equation or inequality.}\\ \msNote{A1}{\(0.7\), or \(\cos 3t\) or \(\sin 3t = \frac{7}{13}\), seen; the inequality is not needed.}} \\
& {\(1.3\cos 3t = 0.7\)\\[4pt] \(t = \frac13\arccos\frac{7}{13} \approx 0.334\)} & M1 A1 & {\msNote{M1}{Equates \(1.3\cos\omega t\) or \(1.3\sin\omega t\), with their \(\omega\), to their \(0.7\), or uses their value of \(\cos 3t\) or \(\sin 3t\), and solves for \(t\).}\\ \msNote{A1}{\(\frac13\arccos\frac{7}{13}\) or awrt 0.334; or the time from \(O\), \(\frac13\arcsin\frac{7}{13}\) or awrt 0.190.}} \\
& {\(T = 4 \times \their{\tfrac13\arccos\tfrac{7}{13}}\)\\[4pt] \(\phantom{T} = \frac43\arccos\frac{7}{13} \approx 1.34\)} & M1 A1 & {\msNote{M1}{Multiplies their time from an extreme position by 4, or subtracts 4 times their time from \(O\) from their period.}\\ \msNote{A1}{Exact form oe, or awrt 1.34.}} \\
\nopagebreak & & [6] & \\
\nopagebreak \SetCell[c=2]{r} Total for Question 3 & & 14 & \\
\end{msscheme}

\msQuestion{4}
\begin{msscheme}{
  hline{3-8,10-13} = {2-Z}{0.4pt, msGrid},
  hline{9,14} = {0.4pt, msGrid},
  row{14} = {font=\bfseries},
}
Question & Scheme & Marks & Notes \\
4(a) & {\((\text{centre} =)\ 7.4, \qquad a = 3.8\)} & B1 & {Both values.} \\
& {\(\tfrac12 T = 6.25 \Rightarrow \omega = \frac{2\pi}{12.5} = \ldots\)} & M1 & {Uses 6 hours 15 minutes as half a period in \(T = \frac{2\pi}{\omega}\). Condone working in minutes.} \\
& {\(\omega = \frac{4\pi}{25}\)} & A1 & {oe, or awrt 0.503 rad per hour, or the equivalent in minutes.} \\
& {\(1.5^2 = \left(\their{\tfrac{4\pi}{25}}\right)^2\left(\their{3.8}^2 - x^2\right)\)\\[4pt] \(\Rightarrow x = \ldots\)} & M1 & {Substitutes \(v = 1.5\), their \(a\) and their \(\omega\) into \(v^2 = \omega^2(a^2 - x^2)\), or uses velocity from \(x = a\cos\omega t\), and solves for \(x\). Units consistent.} \\
& {\(x = \pm 2.35\)} & A1 & {awrt 2.35, or \(x^2\) awrt 5.53; sign not needed.} \\
& {\(7.4 \pm 2.35 \Rightarrow 9.75\,\mathrm{m}\) and \(5.05\,\mathrm{m}\)} & A1 & {Both, awrt 9.75 and 5.05.} \\
\nopagebreak & & [6] & \\
4(b) & {\(3.8\cos\left(\their{\tfrac{4\pi}{25}}\,t\right) = 5.5 - 7.4\)} & M1 & {Sets a cosine model from high tide, or an equivalent model, equal to depth 5.5, using their centre, amplitude and \(\omega\). Condone a sign slip.} \\
& {\(\cos\left(\tfrac{4\pi t}{25}\right) = -\tfrac12\)} & A1 & {Or one correct boundary time, \(t = \frac{25}{6}\) or \(\frac{25}{3}\) hours after 05:40, oe.} \\
& {\(\tfrac{4\pi t}{25} = \tfrac{2\pi}{3}, \tfrac{4\pi}{3} \Rightarrow t = \tfrac{25}{6}, \tfrac{25}{3}\)\\[4pt] \(\tfrac{25}{3} - \tfrac{25}{6} = \ldots\)} & dM1 & {Finds two boundary times within one period from their equation and subtracts, or doubles the time from one boundary to low tide.} \\
& {\(\tfrac{25}{6}\) hours (4 hours 10 minutes)} & A1 & {Or awrt 4.17 hours, or 250 minutes.} \\
\nopagebreak & & [4] & \\
\nopagebreak \SetCell[c=2]{r} Total for Question 4 & & 10 & \\
\end{msscheme}

\msQuestion{5}
\begin{msscheme}{
  hline{3-6,8-12} = {2-Z}{0.4pt, msGrid},
  hline{7,13} = {0.4pt, msGrid},
  row{13} = {font=\bfseries},
}
Question & Scheme & Marks & Notes \\
5(a) & {\(v\frac{\mathrm{d}v}{\mathrm{d}x} = 9(2-x)\)} & M1 & {Uses \(v\frac{\mathrm{d}v}{\mathrm{d}x}\) or \(\frac{\mathrm{d}}{\mathrm{d}x}\left(\frac12 v^2\right)\) for the acceleration in the given equation; or identifies the centre \(x = 2\) and \(\omega = 3\) from it.} \\
& {\(\int v\,\mathrm{d}v = \int (18-9x)\,\mathrm{d}x\)\\[4pt] \(\tfrac12 v^2 = 18x - \tfrac92 x^2 + c\)} & M1 A1 & {\msNote{M1}{Separates the variables and integrates both sides, the power of at least one term increased by 1; or uses \(v^2 = \omega^2\left(a^2 - (x-2)^2\right)\) with their \(\omega\). Condone sign slips.}\\ \msNote{A1}{Correct, oe; the constant, or the amplitude, need not be seen yet.}} \\
& {\(0 = 90 - \tfrac{225}{2} + c \Rightarrow c = \ldots\)} & dM1 & {Substitutes \(x=5\), \(v=0\) into their integrated equation to find the constant, or uses them as limits; or takes the amplitude as the distance from their centre to \(x = 5\).} \\
& {\(c = \tfrac{45}{2}\)\\[4pt] \(v^2 = 36x - 9x^2 + 45\)\\[4pt] \(\phantom{v^2} = 9(5-x)(x+1)\)} & A1~AG & {Correct constant, or amplitude 3, and equation for \(v^2\), leading without error to the printed result; combined steps accepted.} \\
\nopagebreak & & [5] & \\
5(b) & {\(v^2 = 9\left(9-(x-2)^2\right)\)\\[4pt] \(\frac{\mathrm{d}x}{\mathrm{d}t} = 3\sqrt{9-(x-2)^2}\)\\[4pt] \(\int \frac{1}{\sqrt{9-(x-2)^2}}\,\mathrm{d}x = \int 3\,\mathrm{d}t\)} & M1 & {Completes the square as \(p-(x-q)^2\), takes the positive root and separates the variables; or finds the centre and amplitude from part (a). Condone slips in \(p\), \(q\).} \\
& {\(\arcsin\left(\frac{x-2}{3}\right) = 3t + K\)} & M1 A1 & {\msNote{M1}{Integrates to the form \(\arcsin\left(\frac{x-q}{r}\right) = kt + K\), or writes \(x = q + r\sin(kt+K)\) with their centre and amplitude.}\\ \msNote{A1}{Correct, or \(x = 2 + 3\sin(3t+K)\), oe.}} \\
& {\(K = \arcsin\left(\frac{\frac12-2}{3}\right) = \ldots\)} & dM1 & {Substitutes \(t=0\), \(x=\frac12\) into their equation to find the constant.} \\
& {\(K = -\frac{\pi}{6}\)} & A1 & {On the \(x = 2 + 3\sin(3t+K)\) route, needs \(\dot{x} > 0\) at \(t = 0\) used to reject \(-\frac{5\pi}{6}\).} \\
& {\(\arcsin\left(\frac{x-2}{3}\right) = 3t - \frac{\pi}{6}\)\\[4pt] \(x = 2 + 3\sin\left(3t - \frac{\pi}{6}\right)\)} & A1~AG & {Correct equation with \(K=-\frac{\pi}{6}\), rearranged without error to the printed result.} \\
\nopagebreak & & [6] & \\
\nopagebreak \SetCell[c=2]{r} Total for Question 5 & & 11 & \\
\end{msscheme}

\msQuestion{6}
\begin{msscheme}{
  hline{3-5,7-9,11-12} = {2-Z}{0.4pt, msGrid},
  hline{6,10,13} = {0.4pt, msGrid},
  row{13} = {font=\bfseries},
}
Question & Scheme & Marks & Notes \\
6(a) & {\(R - mg = -m\omega^2 x\)} & M1 & {Forms an equation of motion for the box with reaction, weight and acceleration of magnitude \(\omega^2 x\), or \(\omega^2 a\) at an extreme point. Condone sign slips.} \\
& {Greatest \(R = m\left(g + a\omega^2\right)\)\\[4pt] Least \(R = m\left(g - a\omega^2\right)\)} & A1 & {Both correct.} \\
& {\(3m\left(g + a\omega^2\right) = 7m\left(g - a\omega^2\right)\)\\[4pt] \(\Rightarrow \omega^2 = \frac{2g}{5a}\)} & A1~AG & {Correct equation from the \(\frac73\) ratio, leading without error to the printed result.} \\
\nopagebreak & & [3] & \\
6(b) & {\(R = m\left(g - \frac{2gx}{5a}\right) = 0\)\\[4pt] \(\Rightarrow x = \ldots\)} & M1 & {Sets their reaction from (a) to zero, or the downward acceleration \(\omega^2 x\) equal to \(g\), with \(\omega^2 = \frac{2g}{5a}\), and solves for \(x\).} \\
& {\(x = \frac{5a}{2}\)} & A1 & \\
& {\(\frac{5a}{2} < 3a\), so the box leaves the platform} & E1 & {Compares \(\frac{5a}{2}\) with the amplitude \(3a\), or \(3a\omega^2\) with \(g\), and concludes the box leaves the platform.} \\
\nopagebreak & & [3] & \\
6(c) & {\(v^2 = \frac{2g}{5a}\left(9a^2 - \frac{25a^2}{4}\right)\)} & M1 A1 & {\msNote{M1}{Uses \(v^2 = \omega^2\left(A^2 - x^2\right)\) with \(A = 3a\) and their height from (b). Not for amplitude \(a\).}\\ \msNote{A1}{Correct substituted expression.}} \\
& {\(v = \sqrt{\frac{11ga}{10}}\)} & A1 & {oe, such as \(\sqrt{1.1ga}\).} \\
\nopagebreak & & [3] & \\
\nopagebreak \SetCell[c=2]{r} Total for Question 6 & & 9 & \\
\end{msscheme}

\msQuestion{7}
\begin{msscheme}{
  hline{3-6,8-10,12-17} = {2-Z}{0.4pt, msGrid},
  hline{7,11,18} = {0.4pt, msGrid},
  row{18} = {font=\bfseries},
}
Question & Scheme & Marks & Notes \\
7(a) & {\(\frac{36e_0}{0.9} = 0.4 \times 10 \Rightarrow e_0 = 0.1\)} & B1 & \\
& {\(0.4\ddot{x} = 0.4 \times 10 - \frac{36(0.1 + x)}{0.9}\)} & M1 A1 & {\msNote{M1}{Forms Newton's second law with weight and Hooke's-law tension, the extension in terms of a displacement variable. Condone sign slips.}\\ \msNote{A1}{Correct equation, with displacement measured from any point.}} \\
& {\(\ddot{x} = -100x\), so simple harmonic motion} & A1 & {Needs \(x\) measured from equilibrium, or a substitution reaching this form, and a conclusion of simple harmonic motion.} \\
& {\(\text{period} = \frac{2\pi}{\omega} = \frac{2\pi}{10} = \frac{\pi}{5}\)} & M1 A1~AG & {\msNote{M1}{Uses \(\frac{2\pi}{\omega}\) with their \(\omega\) from an equation of the form \(\ddot{x} = -\omega^2 x\).}\\ \msNote{A1~AG}{\(\omega = 10\) from a correct equation of motion, reaching \(\frac{\pi}{5}\) without error.}} \\
\nopagebreak & & [6] & \\
7(b) & {\(3 = 10^2\left(a^2 - 0.1^2\right) \Rightarrow a = \ldots\)} & M1 & {Uses \(v^2 = \omega^2\left(a^2 - x^2\right)\) with \(v = \sqrt{3}\), their \(\omega\) and their equilibrium extension as \(x\), or an energy equation from \(B\) to the slack point with EPE, GPE and KE. Condone the sign of \(x\).} \\
& {\(a = 0.2\)} & A1 & {Or the extension at \(B\), \(0.3\).} \\
& {\(AB = 0.9 + 0.1 + 0.2 = 1.2\,\mathrm{m}\)} & A1 & \\
\nopagebreak & & [3] & \\
7(c) & {\(0 = 3 - 2 \times 10s \Rightarrow s = \ldots\ [= 0.15]\)} & M1 & {Uses \(v^2 = u^2 + 2\alpha s\), or energy, for the free motion from the slack point, with \(v = 0\) and acceleration \(\pm 10\). Condone sign slips.} \\
& {\(BC = 0.2 + 0.1 + 0.15 = 0.45\,\mathrm{m}\)} & A1 & \\
& {\(0.2\cos 10t = -0.1 \Rightarrow t = \ldots\)} & M1 & {Sets their \(a\cos 10t\), or an equivalent sine form with any phase or time added, equal to the slack displacement and solves for \(t\). Condone the sign of \(x\).} \\
& {\(10t = \frac{2\pi}{3} \Rightarrow t = \frac{\pi}{15}\)} & A1 & \\
& {\(0 = \sqrt{3} - 10t_2 \Rightarrow t_2 = \ldots\)\\[4pt] \(\text{time} = \their{\tfrac{\pi}{15}} + \tfrac{\sqrt{3}}{10}\)} & M1 & {Uses \(v = u + \alpha t\), or equivalent, with \(u = \sqrt{3}\), \(v = 0\) and acceleration \(\pm 10\), and adds their time to the slack point.} \\
& {\(\phantom{\text{time}} = \frac{2\pi + 3\sqrt{3}}{30}\,\mathrm{s}\)} & A1 & {Any exact equivalent, such as \(\frac{\pi}{15} + \frac{\sqrt{3}}{10}\).} \\
\nopagebreak & & [6] & \\
\nopagebreak \SetCell[c=2]{r} Total for Question 7 & & 15 & \\
\end{msscheme}

\msQuestion{8}
\begin{msscheme}{
  hline{3-7,9-13,15-17} = {2-Z}{0.4pt, msGrid},
  hline{8,14,18} = {0.4pt, msGrid},
  row{18} = {font=\bfseries},
}
Question & Scheme & Marks & Notes \\
8(a) & {\(\frac{3mge}{5a} = mg\sin\alpha \Rightarrow e = \ldots\)} & M1 & {Resolves parallel to the plane at \(A\) with a Hooke's law tension and a weight component. Condone sin/cos confusion.} \\
& {\(e = a\)} & A1 & \\
& {\(T = \frac{3mg(a-x)}{5a}\)} & B1 & {oe, such as \(\frac{3mg}{5} - \frac{3mgx}{5a}\).} \\
& {\(m\frac{\mathrm{d}^2x}{\mathrm{d}t^2} = \frac{3mg(a-x)}{5a} - mg\sin\alpha\)} & M1 A1 & {\msNote{M1}{Forms an equation of motion along the plane with a tension in terms of \(x\) and the weight component. Condone sign slips.}\\ \msNote{A1}{Correct equation, need not be simplified.}} \\
& {\(\frac{\mathrm{d}^2x}{\mathrm{d}t^2} + \frac{3gx}{5a} = 0\)} & A1~AG & {Correct equation of motion leading without error to the printed equation; combined simplification accepted.} \\
\nopagebreak & & [6] & \\
8(b) & {\(R\sqrt{\frac{3g}{5a}} = 2\sqrt{\frac{3ag}{5}} \Rightarrow R = \ldots\)} & M1 & {Uses maximum speed \(R\omega\), with \(\omega^2 = \frac{3g}{5a}\), equal to the projection speed, or \(v^2=\omega^2(R^2-x^2)\) at \(x=0\).} \\
& {\(R = 2a\)} & A1 & \\
& {\(\their{2a}\sin\omega t = \their{a} \Rightarrow \omega t = \ldots\)} & M1 & {Sets a displacement of the form \(R\sin\omega t\), with their \(R\), equal to their extension at \(A\), or uses an equivalent cosine form with its phase.} \\
& {\(\omega t = \frac{\pi}{6}\)} & A1 & \\
& {\(t = \frac{\pi}{6}\sqrt{\frac{5a}{3g}}\)} & A1 & {oe, such as \(\frac{\pi}{18}\sqrt{\frac{15a}{g}}\).} \\
\nopagebreak & & [5] & \\
8(c) & {\(v^2 = \frac{3g}{5a}\left(\their{2a}^2 - \their{a}^2\right)\)\\[4pt] \(\phantom{v^2} = \ldots\ [= \tfrac{9ag}{5}]\)} & M1 & {Uses \(v^2=\omega^2(R^2-x^2)\) with their \(R\) and \(x =\) their extension at \(A\), or an energy equation from \(A\) including the EPE at \(A\); one such equation to the top earns both M marks.} \\
& {\(\frac12 m\left(\their{\tfrac{9ag}{5}}\right) = mgd\sin\alpha\)\\[4pt] \(\Rightarrow d = \ldots\ [= \tfrac{3a}{2}]\)} & M1 & {Forms an energy equation with KE and GPE, or uses constant acceleration, once the string is slack. Condone sin/cos confusion.} \\
& {\(a + \tfrac{3a}{2} = \tfrac{5a}{2}\)} & A1 & {oe. Needs the distance from \(A\), not from the slack point.} \\
\nopagebreak & & [3] & \\
\nopagebreak \SetCell[c=2]{r} Total for Question 8 & & 14 & \\
\end{msscheme}

\msQuestion{9}
\begin{msscheme}{
  hline{3-4,6-10,12-14,16-21} = {2-Z}{0.4pt, msGrid},
  hline{5,11,15,22} = {0.4pt, msGrid},
  row{22} = {font=\bfseries},
}
Question & Scheme & Marks & Notes \\
9(a) & {\(\frac{3mg(AE-2L)}{2L} = \frac{2mg(6L-AE)}{4L}\)} & M1 & {Equates two tensions of the form \(\frac{\lambda x}{l}\), with both extensions in terms of one unknown length. Condone slips in the extensions.} \\
& {\(3(AE-2L) = 6L-AE \Rightarrow AE = 3L\)} & A1~AG & {Correct equation in one unknown leading without error to \(AE = 3L\).} \\
\nopagebreak & & [2] & \\
9(b) & {\(T_{AP} = \frac{3mg(L+y)}{2L}\)\\[4pt] \(T_{BP} = \frac{2mg(3L-y)}{4L}\)} & B1 & {Both tensions correct in terms of a displacement from \(E\), measured either way, in any letter.} \\
& {\(m\ddot{y} = \frac{mg(3L-y)}{2L} - \frac{3mg(L+y)}{2L}\)} & M1 A1 & {\msNote{M1}{Forms Newton's second law with mass times acceleration and the difference of their two tensions. Condone sign slips.}\\ \msNote{A1}{Correct equation, signs consistent with their direction of displacement.}} \\
& {\(\ddot{y} = -\frac{2g}{L}y \Rightarrow\) simple harmonic motion, centre \(E\)} & A1 & {Correct equation with a conclusion of simple harmonic motion about \(E\), or where \(y = 0\). Strings staying taut need not be shown.} \\
& {Period \(= \frac{2\pi}{\omega} = \frac{2\pi}{\sqrt{\their{\frac{2g}{L}}}}\)} & M1 & {Uses \(\frac{2\pi}{\omega}\) with their \(\omega\) from an equation of the form \(\ddot{y} = -\omega^2 y\).} \\
& {\(\phantom{\text{Period}} = 2\pi\sqrt{\frac{L}{2g}} = \pi\sqrt{\frac{2L}{g}}\)} & A1~AG & {From a correct \(\omega^2 = \frac{2g}{L}\), reaching the printed period without error.} \\
\nopagebreak & & [6] & \\
9(c) & {\(y = \frac52 L - 3L = -\frac12 L\)} & B1 & {Either sign.} \\
& {\(v^2 = \frac{2g}{L}\left(\frac{4L^2}{9} - \frac{L^2}{4}\right)\)} & M1 & {Uses \(v^2 = \omega^2(a^2 - y^2)\) with their \(\omega^2\), \(a = \frac23 L\) and their displacement, or an equivalent energy equation. Not for quoting the formula alone.} \\
& {\(v = \sqrt{\frac{7gL}{18}}\)} & A1 & {Or \(\frac{\sqrt{14gL}}{6}\) oe, or awrt \(0.624\sqrt{gL}\).} \\
\nopagebreak & & [3] & \\
9(d) & {\(\frac{3mg(L+y)}{2L} = 2mg \Rightarrow y = \ldots\)} & M1 & {Sets their tension in \(AP\) equal to, or greater than, \(2mg\) and solves for the displacement or for \(AP\).} \\
& {\(y = \frac13 L\)} & A1 & {Or \(AP = \frac{10}{3}L\).} \\
& {\(y = -\frac23 L\cos\omega t\)\\[4pt] \(-\frac23 L\cos\omega t = \their{\frac13 L} \Rightarrow \cos\omega t = \ldots\)} & M1 & {Uses a cosine or sine model with amplitude \(\frac23 L\) and their boundary displacement to form an equation in \(\omega t\). Condone sign slips.} \\
& {\(\cos\omega t = -\frac12 \Rightarrow \omega t = \frac{2\pi}{3}\)} & A1 & {Or a correct boundary value of \(\omega t\) for their model.} \\
& {\(\left(\frac{4\pi}{3} - \frac{2\pi}{3}\right)\sqrt{\frac{L}{2g}} = \ldots\)} & dM1 & {Finds the interval in \(\omega t\) in one oscillation for which the tension exceeds \(2mg\), and divides by their \(\omega\).} \\
& {\(\frac{\pi}{3}\sqrt{\frac{2L}{g}}\)} & A1 & {Any exact equivalent, such as \(\frac{2\pi}{3}\sqrt{\frac{L}{2g}}\).} \\
\nopagebreak & & [6] & \\
\nopagebreak \SetCell[c=2]{r} Total for Question 9 & & 17 & \\
\end{msscheme}

\msQuestion{10}
\begin{msscheme}{
  hline{3-5,7-8,10-21} = {2-Z}{0.4pt, msGrid},
  hline{6,9,22} = {0.4pt, msGrid},
  row{22} = {font=\bfseries},
}
Question & Scheme & Marks & Notes \\
10(a) & {\(v^2=\left(\sqrt{gl}\right)^2+2gl\Rightarrow v^2=3gl\)} & M1 A1 & {\msNote{M1}{Uses \(v^2=u^2+2as\) with \(u=\sqrt{gl}\), \(a=g\) and \(s=l\). Condone sign slips.}} \\
& {\(\tfrac12 m(3gl)+mgd=\frac{mgd^2}{2l}\)} & M1 A1 & {\msNote{M1}{Forms an energy equation to the lowest point with KE, GPE and EPE in the extension squared. Condone sign slips.}\\ \msNote{A1}{Correct equation oe. A correct energy equation from \(O\) also earns the two marks for \(v^2\).}} \\
& {\(d^2-2ld-3l^2=0\Rightarrow d=3l\)\\[4pt] Greatest length \(l+3l=4l\)} & A1~AG & {Correct energy equation solved without error to \(d=3l\), and the length \(4l\) stated. The negative root need not be shown.} \\
\nopagebreak & & [5] & \\
10(b) & {\(m\ddot{x}=mg-\frac{mgx}{l}\)} & M1 A1 & {\msNote{M1}{Applies Newton's second law to \(P\) with the weight and a tension of the form \(\frac{kx}{l}\). Condone sign slips.}} \\
& {\(\ddot{x}+\frac{g}{l}x=g\Rightarrow\ddot{x}+\omega^2x=g\)} & A1~AG & {Correct equation of motion leading without error to the printed result, with \(\omega^2=\frac{g}{l}\) identified.} \\
\nopagebreak & & [3] & \\
10(c) & {\(x=y+l\Rightarrow\ddot{y}+\omega^2(y+l)=g\)} & M1 & {Substitutes \(x=y+\frac{g}{\omega^2}\), or \(y+l\), with \(\ddot{x}=\ddot{y}\) into the differential equation.} \\
& {\(\ddot{y}+\omega^2y=0\)} & A1 & \\
& {\(\lambda^2+\omega^2=0\Rightarrow y=A\cos\omega t+B\sin\omega t\)} & M1 & {Obtains \(y\) as \(A\cos\omega t+B\sin\omega t\), or an equivalent general solution; the auxiliary equation may be implied.} \\
& {\(x=l+A\cos\omega t+B\sin\omega t\)} & A1 & {oe, with \(l\) or \(\frac{g}{\omega^2}\) added.} \\
& {\(\dot{x}=-A\omega\sin\omega t+B\omega\cos\omega t\)\\[4pt] \(0=l+A,\quad B\omega=\their{\sqrt{3gl}}\)} & M1 & {Uses \(x=0\) and \(\dot{x}=\) their speed from (a) at \(t=0\) in their general solution and its derivative.} \\
& {\(A=-l\)} & A1 & \\
& {\(B=\sqrt{3}\,l\)} & A1 & {oe, such as \(\sqrt{3gl}\sqrt{\frac{l}{g}}\).} \\
& {\(\sqrt{3}\,l\sin\omega t-l\cos\omega t=R\sin(\omega t-\alpha)\)\\[4pt] \(R\cos\alpha=\sqrt{3}\,l,\ R\sin\alpha=l\Rightarrow R=\ldots,\ \alpha=\ldots\)} & M1 & {Uses their \(A\) and \(B\) to find \(R\) and \(\alpha\) for any single-sinusoid form. Condone the ratio inverted.} \\
& {\(R=2l,\quad\alpha=\frac{\pi}{6}\)\\[4pt] \(x=l+2l\sin\left(\omega t-\frac{\pi}{6}\right)\)} & A1 & {Or the correct values for their form.} \\
& {\(l+2l\sin\left(\omega t-\frac{\pi}{6}\right)=0\)\\[4pt] \(\phantom{l}\Rightarrow\sin\left(\omega t-\frac{\pi}{6}\right)=-\frac12\Rightarrow\omega t=\ldots\)} & M1 & {Sets their \(x=0\) and solves for \(\omega t\), from their single sinusoid or by any valid method.} \\
& {\(\omega t-\frac{\pi}{6}=\frac{7\pi}{6}\)} & A1 & {Or \(\omega t=\frac{4\pi}{3}\). Needs the first root after \(t=0\), not \(-\frac{\pi}{6}\).} \\
& {\(t=\frac{4\pi}{3}\sqrt{\frac{l}{g}}\)} & A1 & {Any exact equivalent.} \\
\nopagebreak & & [12] & \\
\nopagebreak \SetCell[c=2]{r} Total for Question 10 & & 20 & \\
\end{msscheme}

\end{document}
